% RECOVERED_RESULT from II REV09. See MAPA_PROCEDENCIA.json.
% Selection of lines and explicit mechanical changes; no new proof.
\subsection{State normalization and the meaning of the modular Hamiltonian}
The second observable is made explicit through the state defining it.
The preceding tensor decomposition is retained: eighteen qutrits
per sector, each with number operator $\operatorname{diag}(0,1,2)$;
$N_m$ is the sum of these operators, and the sectors act on distinct
factors. Let
\[
 z_m=1+e^{-m\Delta_{\rm mod}}+e^{-2m\Delta_{\rm mod}},
 \qquad m\in\{90,120\}.
\]
The state collection in the source has the product realization
\[
 \rho_A=Z^{-1}e^{-\Delta_{\rm mod}(90N_{90}+120N_{120})},
 \qquad Z=z_{90}^{18}z_{120}^{18}.
\]
Commutativity of the number operators and tensor factorization
prove the expression for $Z$. The state has full rank for
real, finite $\Delta_{\rm mod}$, and its spectral logarithm gives
\[
 K_A=-\log\rho_A=(\log Z)I+\Delta_{\rm mod}D.
\]
For $\Delta_{\rm mod}\ne0$, $K_A$ weights the excitations of the two
sectors differently. The joint readout with $N$ recovers both occupations
by equation~\eqref{v:ant:eq:recover}.

The need for the joint readout is apparent in two complementary examples.
The occupation states $(1,0)$ and $(0,1)$ have the same total $N=1$
and different modular weights. Conversely, $(4,0)$ and $(0,3)$ have
the same $D=360$ and different totals. The pair $(N,D)$ separates both
cases. This recovery concerns sector occupations, rather than
each microscopic vector within a degeneracy of those occupations.

When $\Delta_{\rm mod}=0$, the state is proportional to the identity
and $K_A$ is scalar: this readout ceases to distinguish the sectors.
For $\Delta_{\rm mod}\ne0$, the known normalization permits recovery of
$D$ and its composition with area. This defines the role of the modular
Hamiltonian: it characterizes the reduced state and its information distribution.
Its inclusion in the article is motivated by this mathematical function,
whereas the evolution Hamiltonian belongs to the dynamical development
of action.
